% CHAPTER 2: LITERATURE SURVEY
\chapter{Literature Survey}

A comprehensive review of existing literature was conducted to understand the current state of breath analysis technology, IoT-based health monitoring systems, and AI-driven disease prediction methodologies. The following subsections present a critical analysis of seven significant research areas that form the foundation for the proposed system.

\section{VOC-Based Disease Detection}
Studies have established that volatile organic compounds (VOCs) present in exhaled human breath can serve as reliable biomarkers for non-invasive disease detection. Amann et al. (2014) published a comprehensive review in the Journal of Breath Research documenting over 1,800 VOCs identified in human breath, with specific compounds linked to metabolic disorders, respiratory diseases, and inflammatory conditions. Their work demonstrated that acetone concentrations above 1.8 ppm correlate strongly with uncontrolled diabetes mellitus, while elevated levels of pentane and ethane indicate oxidative stress associated with chronic lung diseases. However, most clinical studies used expensive laboratory-grade mass spectrometry equipment (GC-MS) for VOC analysis, making the approach inaccessible for point-of-care applications. The proposed RespiraTech system will address this limitation by utilizing affordable semiconductor gas sensors to detect VOC patterns, making breath analysis accessible for everyday health monitoring.

\section{IoT-Based Electronic Nose Systems}
IoT-based electronic nose (e-nose) systems have emerged as a promising approach for real-time breath analysis. Farraia et al. (2019) reviewed the application of electronic noses in respiratory disease diagnosis, reporting that sensor arrays combined with pattern recognition algorithms could distinguish between healthy subjects and patients with asthma, COPD, and lung cancer with accuracies exceeding 80\%. These systems typically employ metal-oxide semiconductor (MOS) sensors similar to the MQ-series sensors proposed in this project. The key advantage of e-nose systems is their ability to capture the overall ``smell fingerprint'' of a breath sample rather than identifying individual compounds, simplifying the detection pipeline. However, existing implementations often lacked cloud connectivity and real-time data transmission capabilities, which the proposed system will provide through the ESP32's integrated Wi-Fi module.

\section{Machine Learning for Breath Analysis}
Machine learning models have been shown to significantly improve accuracy in disease prediction from breath sensor data compared to threshold-based approaches. Nakhleh et al. (2017) demonstrated in their landmark study published in ACS Nano that an array of gold nanoparticle sensors combined with machine learning classifiers could distinguish between 17 different diseases from breath samples with an average accuracy of 86\%. They employed Random Forest, Discriminant Function Analysis, and Support Vector Machine classifiers, finding that ensemble methods generally outperformed single classifiers. The proposed project will implement similar machine learning approaches using Random Forest and SVM algorithms applied to data from MQ-series gas sensors, adapting these proven classification techniques for a more affordable sensor platform.

\section{AI-Based Respiratory Disease Detection}
AI-based breath analyzers have shown particular effectiveness in detecting respiratory diseases such as asthma and COPD. Dragonieri et al. (2020) used an electronic nose with eight MOS sensors to distinguish between patients with mild, moderate, and severe asthma with a classification accuracy of 89\% using principal component analysis and neural network classifiers. Their work demonstrated that sensor-derived breath profiles can capture subtle differences in airway inflammation levels. However, their system was clinic-based and did not support remote monitoring. The proposed RespiraTech system will extend this capability by enabling continuous, remote breath monitoring through cloud-based data processing and mobile application delivery.

\section{Breath-Based Diabetes Screening}
Breath analysis has shown significant promise for diabetes detection by identifying elevated acetone levels in exhaled air. Righettoni et al. (2010) developed a Si-doped WO$_3$ nanoparticle sensor specifically designed for breath acetone detection, achieving a detection limit of 20 ppb. Their research confirmed that breath acetone concentrations in diabetic patients (typically 1.8--5.0 ppm) are significantly higher than in healthy individuals (0.3--0.9 ppm). The MQ-3 sensor included in the proposed system is sensitive to ethanol and acetone vapours, enabling the detection of diabetes-related breath biomarkers. While the MQ-3 lacks the specificity of purpose-built nanoparticle sensors, the use of multi-sensor fusion and machine learning pattern recognition will compensate for individual sensor limitations.

\section{Multi-Sensor Fusion Approaches}
Multi-sensor systems have been shown to increase reliability and precision of breath analysis compared to single-sensor configurations. Binson et al. (2021) demonstrated that an array of four different gas sensors combined with a neural network classifier achieved 94\% accuracy in distinguishing between healthy and diabetic breath samples, compared to only 72\% accuracy when using a single sensor. Their sensor fusion approach leveraged the complementary sensitivities of different sensor types to capture a broader spectrum of VOC information. The proposed RespiraTech system will adopt this multi-sensor strategy by combining MQ-135 (broad-spectrum gas detection) with MQ-3 (alcohol/acetone-specific) and DHT11/DHT22 (environmental compensation), creating a comprehensive sensor array for reliable breath profiling.

\section{Cloud-Based IoT Health Monitoring}
Cloud-based IoT platforms have been shown to enable effective real-time health monitoring and data storage. Dey et al. (2018) developed a cloud-connected IoT health monitoring system using ThingSpeak for data aggregation and visualization, demonstrating that MQTT and HTTP-based data transmission protocols can deliver sensor readings to cloud platforms with latency under 2 seconds over standard Wi-Fi connections. Their work established the feasibility of real-time cloud-IoT integration for health applications. The proposed system will leverage similar cloud platforms (ThingSpeak or Firebase) for data storage, ML model deployment, and result delivery through REST APIs to the mobile/web application.

\section{Research Gap}
The literature review reveals that while significant progress has been made in breath-based disease detection, existing solutions suffer from one or more critical limitations:

\begin{enumerate}[leftmargin=1.5cm]
\item \textbf{High Cost:} Most clinical breath analysis systems use expensive laboratory instruments (GC-MS, SIFT-MS) that are inaccessible for routine screening and point-of-care applications.

\item \textbf{Lack of IoT Integration:} Many research prototypes are standalone devices without cloud connectivity, real-time data transmission, or remote monitoring capabilities.

\item \textbf{Single-Disease Focus:} Most systems target a single disease (e.g., only diabetes or only asthma) rather than providing multi-disease screening from a single breath sample.

\item \textbf{No User-Friendly Interface:} Research prototypes typically display raw sensor values without providing interpreted health risk assessments, percentages, or actionable alerts for non-medical users.
\end{enumerate}

This project will address all four limitations by combining affordable MQ-series gas sensors with ESP32-based IoT connectivity, cloud-based multi-disease AI prediction, and a user-friendly mobile/web application for health risk visualization and alerting.
